\begin{textsample}{POS Dim 1 – rn – Score 40.00 – rn/rn\_em\_43.txt}  \label{ex:f1_pos_267}
Independentemente de qualquer estratégia selecionada para determinado contexto escolar, estas devem enfatizar a formação integral, destacando o desenvolvimento dos procedimentos e a construção de valores por meio do uso de diferentes linguagens.
O estudo das \textbf{ciências} da natureza é uma importante fonte de evidências para fundamentar as tomadas de decisão conscientes ao longo da vida.
Munidos de uma \textbf{compreensão} ampla do conhecimento, com a habilidade de \textbf{sistematizar} as \textbf{inúmeras} informações disponíveis, os estudantes têm a possibilidade de se preparar para o cenário de constantes mudanças do \textbf{século} XXI e de construir seu projeto de vida a partir do desenvolvimento de suas \textbf{capacidades} de resolução de \textbf{problemas}, pensando e agindo de \textbf{maneira} criativa, capazes de modificar a realidade com responsabilidade social, colaboração, empatia e autocuidado.
Ao considerarmos o contexto \textbf{atual} da informação, em que os computadores do mundo todo estão conectados em rede a qualquer hora, se faz necessário ressignificar as relações escolares entre professores, estudantes e objetos de conhecimento ( DAVIDSON, 2012 ).
A quantidade de informação acumulada até aqui pela \textbf{humanidade} é imensurável, tornando difícil a retenção por um indivíduo de \textbf{tantas} informações disponíveis e conectadas.
Por isso, o professor passa a ter um papel de curador do conhecimento, conduzindo os estudantes às fontes de informações \textbf{relevantes} e confiáveis para seus projetos, catalisando o processo de desenvolvimento por meio de intervenções que levem à reflexão e facilitem, assim, as ações do estudante na construção compartilhada do conhecimento.
Entretanto, não devemos fazer uso da \textbf{tecnologia} \textbf{digital} apenas como entretenimento, mas usá -la como uma ferramenta para o aprendizado e criatividade ( RESNICK, 2017 ), permitindo que os estudantes analisem \textbf{problemas}, extraiam dados e se expressem por meio das \textbf{inúmeras} estruturas tecnológicas disponíveis nos tempos \textbf{atuais}.
Um dos fatores importantes a ser considerado pelos docentes nas \textbf{abordagens} de \textbf{Ciências} da Natureza para o \textbf{aluno} que chega ao ensino \textbf{médio} é o seu percurso que se iniciou nas etapas anteriores.
É \textbf{imprescindível} que a equipe docente conheça os níveis de \textbf{compreensão} dos estudantes, seja no âmbito cognitivo, seja em suas atitudes e valores desenvolvidos nos anos escolares antecedentes.
Como um espiralamento, \textbf{competências} e habilidades do Ensino Fundamental serão revisitadas no Ensino Médio, seguindo um percurso de \textbf{aprofundamento}.
Por exemplo, no Ensino Fundamental, o estudante deve desenvolver a habilidade “ Propor iniciativas individuais e coletivas para a solução de \textbf{problemas} ambientais da cidade ou da comunidade, com base na \textbf{análise} de ações de consumo consciente e de sustentabilidade bem sucedidas - EF09CI13 ” ( SEEC, 2018 ), que por sua vez está relacionada à \textbf{competência} específica “ Avaliar aplicações e implicações políticas, \textbf{socioambientais} e culturais da \textbf{ciência} e de suas \textbf{tecnologias} para propor alternativas aos desafios do mundo \textbf{contemporâneo}, incluindo aqueles relativos ao mundo do trabalho - EFCNC4 ” ( SEEC, 2018 ).
Já no ensino \textbf{médio}, essas habilidades e \textbf{competências} serão retomadas de alguma \textbf{maneira} na habilidade “ Justificar a importância da preservação e conservação da biodiversidade, considerando \textbf{parâmetros} \textbf{qualitativos} e quantitativos, e avaliar os efeitos da ação humana e das políticas ambientais para a garantia da sustentabilidade do planeta - EM13CNT206 ” ( BNCC, 2018 ).
Na prática, em um determinado projeto, enquanto os estudantes no ensino fundamental precisarão propor iniciativas para soluções de \textbf{problemas} ambientais - como o manejo do lixo e o destino do óleo usado nos restaurantes locais ; no ensino \textbf{médio}, eles precisarão se posicionar, justificar e \textbf{argumentar}.
Ainda, por exemplo, ter que justificar a relevância da proteção dos principais destinos turísticos do litoral potiguar, diante do empreendedorismo hoteleiro e gastronômico neles instalado e a preservação das nossas matas e biomas, como nas Áreas de Proteção Ambiental ( APA ) do RN.
Para isso, pressupõe -se que eles necessitarão de habilidades adquiridas anteriormente no ensino fundamental, como avaliar situações e propor alternativas.
Todas essas habilidades juntas se relacionam com a \textbf{competência} geral de pensar criticamente e cientificamente, sendo esse um importante objetivo a ser alcançado na formação dos cidadãos na educação básica.
Desse modo, a escola potiguar deve proporcionar uma formação que \textbf{auxilie} a transição do jovem do Rio Grande do Norte ao mundo adulto, considerando sua \textbf{pluralidade}, visando a construção e reconstrução permanente da sua identidade individual e coletiva.
Deve estar \textbf{focada} na formação do cidadão em todos os papéis sociais, considerando a construção do projeto de vida dos jovens, não se limitando à instrumentalização do estudante para o mundo do trabalho.
Nesse sentido, as Ciências da Natureza contribuem para esse processo, \textbf{fomentando} o pensamento \textbf{crítico} \textbf{baseado} em evidências científicas, munindo o estudante de condições de se tornar um \textbf{aprendiz} permanente, além de permitir que seja capaz de se apropriar das \textbf{tecnologias} \textbf{emergentes}, assim como refletir sobre seu uso e impacto na vida coletiva.
O Referencial Curricular do Ensino Médio do Rio Grande do Norte é organizado em \textbf{matriz} de área e \textbf{matrizes} de componentes curriculares.
Na primeira, correlacionam -se as três \textbf{competências} específicas da BNCC, as habilidades da BNCC vinculadas às \textbf{competências} específicas, a articulação da habilidade no componente e os objetos do conhecimento.
Enquanto a segunda é composta por eixos integradores, habilidades específicas, objetos de conhecimento e sugestões didáticas.
Parte -se da \textbf{análise} geral da \textbf{matriz} da área, abrangendo todos os componentes curriculares para questões mais circunscritas, onde se estudam as \textbf{teorias} e os modelos mais específicos de Física, Química e Biologia, facilitando assim as contribuições de cada um deles para o estudo das Ciências da Natureza e suas Tecnologias.
O processo \textbf{avaliativo} na Área de Ciências da Natureza e suas Tecnologias enfatiza os aspectos de caráter democrático, integral e inclusivo.
É preciso avaliar o domínio dos procedimentos da investigação científica, conduzindo o indivíduo à ampliação da sua \textbf{capacidade} \textbf{dialética}, de interpretação de fenômenos naturais, da apropriação da linguagem específica da área, da argumentação sobre questões científicas, tecnológicas e \textbf{socioambientais} ; e para a utilização dos conhecimentos de forma ética, responsável e socialmente colaborativa.
Para \textbf{tanto}, os professores da Área de Ciências da Natureza podem lançar mão de instrumentos que possibilitem avaliar a \textbf{progressão} da aprendizagem dos estudantes diante da \textbf{compreensão} e desenvolvimento das \textbf{competências} e habilidades da Área e dos eixos integradores, além de habilidades específicas dos componentes, considerando os quatro eixos integradores ( Energia : fontes, transformações e seu uso sustentável ; Modelos científicos x realidade : validade e utilização na sociedade ; Meio ambiente e saúde ; Investigação e divulgação científica e tecnológica )

% matched lemmas: abordagem, aluno, análise, aprendiz, aprofundamento, argumentar, atual, auxiliar, avaliativo, basear, capacidade, ciência, competência, compreensão, contemporâneo, crítico, dialético, digital, emergente, focar, fomentar, humanidade, imprescindível, inúmero, maneira, matriz, médio, parâmetro, pluralidade, problema, progressão, qualitativo, relevante, sistematizar, socioambiental, século, tanto, tecnologia, teoria
\end{textsample}
